% ═══════════════════════════════════════════════════════════════
% MT-02a: Minitest Familienvokabular
% Spanisch Klasse 7
% ═══════════════════════════════════════════════════════════════
% Dauer: 10 Minuten
% Punkte: 26
% Inhalt: Familienmitglieder, Artikel, Lückentext
% ═══════════════════════════════════════════════════════════════

\documentclass[11pt, a4paper]{article}

\usepackage[utf8]{inputenc}
\usepackage[T1]{fontenc}
\usepackage[ngerman]{babel}
\usepackage{helvet}
\renewcommand{\familydefault}{\sfdefault}

\usepackage[margin=2cm]{geometry}
\usepackage{enumitem}
\usepackage{tabularx}
\usepackage{booktabs}

\usepackage{xcolor}
\usepackage{tcolorbox}
\tcbuselibrary{skins}

\usepackage{fancyhdr}
\usepackage{lastpage}
\usepackage{needspace}

% FARBEN
\definecolor{spanisch-color}{HTML}{c41e3a}
\definecolor{grau-color}{HTML}{888888}
\definecolor{hellgrau-color}{HTML}{f5f5f5}

\colorlet{spanisch}{spanisch-color}
\colorlet{grau}{grau-color}
\colorlet{hellgrau}{hellgrau-color}

\newcommand{\fachfarbe}{spanisch}

\newcommand{\thema}{Minitest: Familienvokabular}
\newcommand{\fach}{Spanisch}
\newcommand{\klasse}{7}
\newcommand{\maxpunkte}{26}
\newcommand{\zeit}{10 Minuten}
\newcommand{\datum}{\today}

\pagestyle{fancy}
\fancyhf{}
\renewcommand{\headrulewidth}{0.4pt}
\renewcommand{\footrulewidth}{0.4pt}

\lhead{\fach{} -- Klasse \klasse}
\chead{\textbf{MT-02a}}
\rhead{\datum}

\lfoot{\zeit}
\cfoot{}
\rfoot{Seite \thepage\ von \pageref{LastPage}}

\newcommand{\punktebox}[1]{\hfill\fbox{\small \_/#1}}
\newcommand{\luecke}[1]{\rule{#1}{0.4pt}}

\newtcolorbox{infobox}{
    colback=hellgrau,
    colframe=\fachfarbe,
    boxrule=1pt,
    arc=0pt,
    left=8pt, right=8pt, top=8pt, bottom=8pt
}

\newenvironment{aufgabe}[1]{%
    \needspace{5\baselineskip}%
    \nopagebreak[4]%
    \vspace{10pt}
    \noindent\textbf{\textcolor{\fachfarbe}{Aufgabe #1}}
    \nopagebreak[4]%
    \vspace{5pt}
    \nopagebreak[4]%
    \begin{list}{}{\leftmargin=0pt}
    \item[]
}{%
    \end{list}
}

\newcommand{\es}[1]{\textcolor{\fachfarbe}{\textbf{#1}}}

\begin{document}

% ═══════════════════════════════════════════════════════════════
% KOPFBEREICH
% ═══════════════════════════════════════════════════════════════

\begin{center}
    {\Large\bfseries\textcolor{\fachfarbe}{\thema}}
\end{center}

\vspace{5pt}

\noindent
\begin{tabularx}{\textwidth}{|X|X|X|}
\hline
\textbf{Name:} \luecke{4cm} &
\textbf{Klasse:} \klasse &
\textbf{Datum:} \luecke{2cm} \\
\hline
\end{tabularx}

\vspace{10pt}

\begin{infobox}
\textbf{Zeit:} \zeit{} | \textbf{Punkte:} \maxpunkte{} | \textbf{Hilfsmittel:} keine
\end{infobox}

\vspace{15pt}

% ═══════════════════════════════════════════════════════════════
% AUFGABE 1: Zuordnung Familienmitglieder (8 Punkte)
% ═══════════════════════════════════════════════════════════════

\begin{aufgabe}{1 -- Familienmitglieder zuordnen}
Verbinde die spanischen Familienbegriffe mit der deutschen Bedeutung. \punktebox{8}
\end{aufgabe}

\begin{center}
\begin{tabular}{rcl}
\es{el padre} & \luecke{2cm} & die Schwester \\[0.8em]
\es{la madre} & \luecke{2cm} & der Vater \\[0.8em]
\es{el hermano} & \luecke{2cm} & die Großmutter \\[0.8em]
\es{la hermana} & \luecke{2cm} & die Mutter \\[0.8em]
\es{el abuelo} & \luecke{2cm} & die Tante \\[0.8em]
\es{la abuela} & \luecke{2cm} & der Bruder \\[0.8em]
\es{el tío} & \luecke{2cm} & der Großvater \\[0.8em]
\es{la tía} & \luecke{2cm} & der Onkel \\
\end{tabular}
\end{center}

\vspace{15pt}

% ═══════════════════════════════════════════════════════════════
% AUFGABE 2: Artikel einsetzen (8 Punkte)
% ═══════════════════════════════════════════════════════════════

\begin{aufgabe}{2 -- Artikel einsetzen}
Ergänze den richtigen Artikel: \es{el} oder \es{la}. \punktebox{8}
\end{aufgabe}

\noindent
\begin{tabularx}{\textwidth}{|c|X|X|}
\hline
\textbf{Artikel} & \textbf{Familienmitglied} & \textbf{Bedeutung} \\
\hline
\luecke{1.5cm} & primo & der Cousin \\[0.6em]
\hline
\luecke{1.5cm} & prima & die Cousine \\[0.6em]
\hline
\luecke{1.5cm} & padre & der Vater \\[0.6em]
\hline
\luecke{1.5cm} & madre & die Mutter \\[0.6em]
\hline
\luecke{1.5cm} & hermano & der Bruder \\[0.6em]
\hline
\luecke{1.5cm} & hermana & die Schwester \\[0.6em]
\hline
\luecke{1.5cm} & abuelo & der Großvater \\[0.6em]
\hline
\luecke{1.5cm} & abuela & die Großmutter \\[0.6em]
\hline
\end{tabularx}

\vspace{15pt}

% ═══════════════════════════════════════════════════════════════
% AUFGABE 3: Lückentext (10 Punkte)
% ═══════════════════════════════════════════════════════════════

\begin{aufgabe}{3 -- Lückentext}
Ergänze das passende Familienmitglied (ohne Artikel). \punktebox{10}
\end{aufgabe}

\noindent
a) Mi \luecke{3cm} se llama Carlos. \hfill (mein Vater)

\vspace{0.8em}

\noindent
b) Mi \luecke{3cm} se llama Ana. \hfill (meine Mutter)

\vspace{0.8em}

\noindent
c) Tengo un \luecke{3cm}. \hfill (einen Bruder)

\vspace{0.8em}

\noindent
d) Tengo una \luecke{3cm}. \hfill (eine Schwester)

\vspace{0.8em}

\noindent
e) Mi \luecke{3cm} es muy simpático. \hfill (mein Großvater)

% ═══════════════════════════════════════════════════════════════
% PUNKTEÜBERSICHT
% ═══════════════════════════════════════════════════════════════

\vspace{25pt}
\hrule
\vspace{10pt}

\noindent
\begin{tabularx}{\textwidth}{|X|c|c|c||c|}
\hline
\textbf{Aufgabe} & 1 & 2 & 3 & \textbf{Gesamt} \\
\hline
\textbf{max. Punkte} & 8 & 8 & 10 & \textbf{26} \\
\hline
\textbf{erreicht} & & & & \\[0.8em]
\hline
\end{tabularx}

\vspace{10pt}

\begin{flushright}
\textbf{Note:} \luecke{2cm}
\end{flushright}

\end{document}
